\section{Other directions and the remaining gaps}\label{sec:remaining}

\subsection{An edge-frequency consequence}

The following implication was proposed in the supplied research list.
It follows from the source manuscript's formalized forest estimate.
It is recorded here with its denominator and graph assumptions.

\begin{proposition}\label{prop:research-edge}
Let $G$ be connected and simple, with sufficiently large rank $r$.
Suppose an edge $e$ lies on a cycle $W$ of length at most $\sqrt r$.
Then
\[
 c_e(G)\le\left(\frac14+\frac{10}{\log_2\log_2\log_2 r}\right)2^r.
\]
If also $c(G)\ge(1/2-\eta)2^r$ with $0\le\eta<1/2$, then
\[
 \frac{c_e(G)}{c(G)}
 \le\frac{1/4+10/(\log_2\log_2\log_2 r)}{1/2-\eta}.
\]
\end{proposition}

\begin{proof}
Restriction outside $W$ has fibers of size two, namely $\{x,x\mathbin\triangle E(W)\}$.
Exactly one member of each fiber contains $e$.
A cycle containing $e$ has a forest outside part.
Thus $c_e(G)\le p(G,W)2^{r-1}$.
Apply the source manuscript's forest theorem and divide by the stated lower bound on $c(G)$.
\end{proof}

A connected simple graph of minimum degree three has at most $2r-2$ vertices.
A breadth-first search gives a cycle of length $O(\log r)$.
Hence such a graph has an edge satisfying the proposition for large $r$.
For families with $c(G)/2^r\to1/2$, this gives an edge in at most $(1/2+o(1))c(G)$ cycles.
This is not the exact edge-half theorem for arbitrary graphs.

At a vertex $v$ of degree $d\ge4$, the exact local half bound is elementary:
\[
 \sum_{e\ni v}c_e(G)=2c_v(G)\le2c(G).
\]
Some incident edge therefore has $c_e(G)\le c(G)/2$.
A bridge lies in no cycle.
Thus a global counterexample with minimum degree at least three would have to be cubic and bridgeless.

\subsection{Why the general error-term questions remain}

The extended marked-set range in Section~\ref{sec:research-path} concerns a fixed positive error.
It does not supply polynomial decay in the error as $r$ grows.
The source proof's weighted cycle supply remains logarithmic in its available length ratio.
The moment-only example in Section~\ref{sec:inverse} prevents stronger conclusions from those moments alone.

The chord construction gives
\[
 \frac{c(G)}{2^r}\le p(G^+,W)+2^{-t}.
\]
Its witness can have $t$ components.
The current forest estimate allows only $O(\log r)$ components.
Consequently, this direct use of the construction cannot make $2^{-t}$ exponentially small in $r$.
An exponential excess bound needs an additional argument for the omitted cycles or a different witness.

The fixed-edge-budget problem needs a different improvement as well.
Multiplying $2^r$ by a constant does not improve its exponential base when $r$ is proportional to the edge count.
No result here proves the proposed base $(2+2\sqrt2)^{1/5}$.
The degree-sequence paper supplying that target is Dvo\v r\'ak, Morrison, Noel, Norin, and Postle,
\emph{Bounding the number of cycles in a graph in terms of its degree sequence},
\url{https://arxiv.org/abs/1907.12091}.

\subsection{A boundary for matroid extensions}

Let $M=M^*(K_{r+1})$ be the cographic matroid of the complete graph.
Its binary cycle space is the graph's cut space and has dimension $r$.
Every nonempty cut $\delta(S)$ is a bond, since both induced sides are connected.
Thus every nonzero vector is a circuit, and $M$ has $2^r-1$ circuits.
The one-half theorem cannot extend to all binary matroids or even all regular matroids.
This counterexample was already identified in the supplied research list.

\subsection{Targets not resolved by this work}

General linear path stability remains open in this note.
The exact four-regular bound also remains open beyond the reductions and finite checks.
The random-host theorem leaves both the sharp correction constant and the full-threshold upper bound open.
No new result here proves monotonicity of minimum pancyclic excess or the proposed edge-pancyclic density bound.
The supplied list's distinction between cycle-space counts and obligations through every edge remains essential.
